\section{Part 1 results}

\subsection{Learning curves and crossover}

\Cref{fig:learning-curves} shows all nine pipelines. TF-IDF SVM or logistic regression leads at small and medium sizes. Neural curves rise more steeply at large budgets, although the threshold differs by domain.

\begin{figure*}[t]
  \centering
  \includegraphics[width=\textwidth]{part1_learning_curves.pdf}
  \caption{Part~1 test macro-F1. Points are three-seed means and error bars are $\pm$ one sample SD. The x-axis is logarithmic and labels show actual tested counts.}
  \label{fig:learning-curves}
\end{figure*}

\begin{table}[t]
\centering
\caption{Part 1 crossover summary. Sustained means the $\geq 0.01$ mean macro-F1 margin holds at every larger tested size.}
\label{tab:crossover}
\small
\resizebox{\columnwidth}{!}{%
\begin{tabular}{lrr}
\toprule
Dataset & First observed & First sustained \\
\midrule
NSMC & 64,000 & 64,000 \\
Shopping & Not through 159,918 & Not through 159,918 \\
Steam & Not through 79,828 & Not through 79,828 \\
\bottomrule
\end{tabular}
}
\end{table}


NSMC crosses at \Result{part1_crossover_nsmc} labels and stays above the criterion at full data. Shopping and Steam never reach a sustained one-point advantage within their grids. The full-data neural gaps are \Result{story_full_gap_nsmc} points on NSMC, \Result{story_full_gap_shopping} on Shopping, and \Result{story_full_gap_steam} on Steam. The latter two are positive but smaller than the prespecified threshold.

\begin{figure*}[t]
  \centering
  \includegraphics[width=0.98\textwidth]{part1_crossover.pdf}
  \caption{Difference between the best complete neural and TF-IDF mean at each size. Green points meet the 0.01 criterion; the dashed line marks the threshold.}
  \label{fig:crossover}
\end{figure*}

\subsection{Full-data ranking}

\begin{table*}[t]
\centering
\caption{Full-data test macro-F1 (percentage points), mean $\pm$ sample SD ($n=3$). Bold marks the unique best rounded mean; underline marks the unique second best.}
\label{tab:full-data}
\scriptsize
\begin{tabular}{lccc}
\toprule
Method & NSMC & Shopping & Steam \\
\midrule
TF-IDF + LR & 80.07 $\pm$ 0.00 (3) & 88.34 $\pm$ 0.00 (3) & 75.82 $\pm$ 0.00 (3) \\
TF-IDF + SVM & 80.04 $\pm$ 0.00 (3) & 88.30 $\pm$ 0.00 (3) & 75.67 $\pm$ 0.00 (3) \\
TF-IDF + RF & 75.39 $\pm$ 0.01 (3) & 85.09 $\pm$ 0.02 (3) & 72.48 $\pm$ 0.04 (3) \\
TF-IDF + XGBoost & 68.94 $\pm$ 0.03 (3) & 80.37 $\pm$ 0.05 (3) & 68.21 $\pm$ 0.09 (3) \\
Avg. embeddings + MLP & 81.36 $\pm$ 0.03 (3) & 88.62 $\pm$ 0.19 (3) & 76.39 $\pm$ 0.14 (3) \\
LSTM & 81.64 $\pm$ 0.36 (3) & 88.97 $\pm$ 0.13 (3) & \textbf{76.68 $\pm$ 0.44 (3)} \\
GRU & \underline{81.80 $\pm$ 0.02 (3)} & \underline{89.04 $\pm$ 0.21 (3)} & 76.60 $\pm$ 0.42 (3) \\
Transformer & \textbf{81.85 $\pm$ 0.02 (3)} & \textbf{89.14 $\pm$ 0.11 (3)} & \underline{76.64 $\pm$ 0.21 (3)} \\
Mamba & 80.72 $\pm$ 0.12 (3) & 88.09 $\pm$ 0.08 (3) & 75.39 $\pm$ 0.14 (3) \\
\bottomrule
\end{tabular}
\end{table*}


Transformer is the full-data leader on NSMC and Shopping, with mean macro-F1 of \Result{story_full_best_neural_nsmc} and \Result{story_full_best_neural_shopping} percentage points. LSTM leads Steam at \Result{story_full_best_neural_steam}. The small differences among the leading neural models should not be read as formal rankings beyond these runs.

\begin{figure*}[t]
  \centering
  \includegraphics[width=0.98\textwidth]{full_data_comparison.pdf}
  \caption{Full-data macro-F1 with sample-SD error bars. Axes share the same range; no truncation is used to magnify within-domain gaps.}
  \label{fig:full-data}
\end{figure*}

\subsection{Efficiency}

\begin{table*}[t]
\centering
\caption{Representative full-data efficiency. Times are means across three seeds on one recorded hardware ID; TF-IDF used CPU and neural methods used CUDA, so cross-device latency is descriptive.}
\label{tab:efficiency}
\scriptsize
\begin{tabular}{llrrrrr}
\toprule
Dataset & Method & Macro-F1 (\%) & Fit (s) & End-to-end p50 (ms) & Size (MB) & Peak memory (MB) \\
\midrule
NSMC & TF-IDF + LR & 80.07 & 5.1 & 0.50 & 0.42 & 1702 \\
NSMC & Transformer & 81.85 & 235.9 & 1.87 & 11.23 & 335 \\
Shopping & TF-IDF + LR & 88.34 & 5.9 & 0.50 & 0.43 & 3483 \\
Shopping & Transformer & 89.14 & 289.7 & 1.96 & 11.24 & 333 \\
Steam & TF-IDF + LR & 75.82 & 3.4 & 0.52 & 0.42 & 2162 \\
Steam & LSTM & 76.68 & 95.1 & 1.03 & 12.13 & 176 \\
\bottomrule
\end{tabular}
\end{table*}


The strongest neural method used \Result{story_fit_ratio_nsmc} times the mean fit time of the strongest TF-IDF method on NSMC, \Result{story_fit_ratio_shopping} times on Shopping, and \Result{story_fit_ratio_steam} times on Steam. Neural artifacts were also much larger. In return, full-data macro-F1 improved modestly, and only NSMC cleared the one-point criterion. Linear TF-IDF offers the clearest low-cost choice unless the last fraction of performance is worth GPU training and a larger deployed model.

\begin{figure*}[t]
  \centering
  \includegraphics[width=0.98\textwidth]{efficiency_pareto.pdf}
  \caption{Full-data quality versus mean fit time. Outlined points are Pareto-efficient within a dataset. Fit time is plotted on a log scale.}
  \label{fig:pareto}
\end{figure*}

The confusion-matrix check in \Cref{fig:confusion} shows that the NSMC neural gain does not come from sacrificing one label; error rates fall for the positive class while remaining comparable for the negative class.
\FloatBarrier
